\documentclass{article}
\usepackage[margin=28mm]{geometry}
\usepackage{amsmath,amssymb}
\usepackage{xurl}
\usepackage[hidelinks]{hyperref}
% Shared commands for notes in docs/paper-gaps/.

\DeclareUnicodeCharacter{2080}{\ifmmode{}_{0}\else\textsubscript{0}\fi}
\DeclareUnicodeCharacter{2081}{\ifmmode{}_{1}\else\textsubscript{1}\fi}
\DeclareUnicodeCharacter{2082}{\ifmmode{}_{2}\else\textsubscript{2}\fi}
\DeclareUnicodeCharacter{2099}{\ifmmode{}_{n}\else\textsubscript{n}\fi}

\newcommand{\Lean}{\textsc{Lean}}
\ExplSyntaxOn
\NewDocumentCommand{\leanid}{m}{
  \ifmmode
    \textsf{\detokenize{#1}}
  \else
    \group_begin:
    \urlstyle{sf}
    \tl_set:Nn \l_tmpa_tl {#1}
    \tl_replace_all:Nnn \l_tmpa_tl {\_} {_}
    \exp_args:NV \path \l_tmpa_tl
    \group_end:
  \fi
}
\ExplSyntaxOff
\providecommand{\tr}{\operatorname{tr}}
\newcommand{\tnleanIssueBase}{https://github.com/LionSR/TNLean/issues/}
\newcommand{\tnleanPrBase}{https://github.com/LionSR/TNLean/pull/}
\newcommand{\tnleanDocBase}{https://sirui-lu.com/TNLean/docs/}
\newcommand{\ghissue}[1]{\href{\tnleanIssueBase#1}{GitHub issue \##1}}
\newcommand{\ghpr}[1]{\href{\tnleanPrBase#1}{PR \##1}}
\newcommand{\doclink}[1]{\href{\tnleanDocBase#1.html}{\path{#1}}}

% Dirac notation and the identity operator, as in blueprint/src/macros/common.tex.
% \providecommand keeps note-local definitions authoritative.
\usepackage{dsfont}
\providecommand{\ket}[1]{|#1\rangle}
\providecommand{\bra}[1]{\langle #1|}
\providecommand{\braket}[2]{\langle #1 | #2 \rangle}
\providecommand{\ketbra}[2]{|#1\rangle\!\langle #2|}
\providecommand{\Id}{\mathds{1}}
\providecommand{\one}{\mathds{1}}
\providecommand{\C}{\mathbb{C}}
\providecommand{\MN}[1]{M_{#1}(\C)}
\providecommand{\MD}{M_D(\C)}

% Verdict marker: \gapnote{<kind>}{<status>}. Kind is the policy
% classification (clarification, local-correction, scope-restriction,
% unfaithful, false-source, open-gap); status is open, wip (elimination
% underway), resolved, or historical. Machine-read by the site index;
% prints a verdict line.
\newcommand{\gapnote}[2]{\par\noindent\textbf{Verdict:} #1 (#2).\par}

\providecommand{\gapnote}[2]{\par\noindent\textbf{Verdict:} #1 (#2).\par}
\title{The Number of Final States in Hastings' Hierarchical Construction}
\author{}
\date{}
\begin{document}
\maketitle
\gapnote{clarification}{resolved}

\section{The asserted count}
Hastings, \emph{Solving gapped Hamiltonians locally}~\cite{Hastings2006Solving},
\href{https://arxiv.org/abs/cond-mat/0508554v3}{arXiv:cond-mat/0508554v3},
source \path{proj3new.tex}, lines 722--730, ends the hierarchical construction
of Section~II with the assertion that, once every final state $\alpha^h$ lies
within energy $\Delta E/2$ of the ground state, ``it can then be shown that
the number of distinct $\alpha^h$ is at most $2n$, where $n$ is the number of
ground states''. No argument is given. This note supplies one and records the
form in which the hypothesis is used.

\section{The argument}
Let $P$ be the orthogonal projection onto the ground states, so that
$\operatorname{tr}P=n$, and suppose that the Hamiltonian $H$ has gap $\Delta E$
above the ground energy $E_0$ in the sense
\begin{align}
 H-E_0\ge\Delta E\,(1-P).
 \label{eq:hastings05_final_state_count_gap}
\end{align}
Let $\phi_1,\dots,\phi_K$ be orthonormal with
$\langle\phi_k,H\phi_k\rangle\le E_0+\theta\Delta E$ for some $\theta<1$. By
(\ref{eq:hastings05_final_state_count_gap}),
$\Delta E\,(1-\langle\phi_k,P\phi_k\rangle)\le\theta\Delta E$, so
$\langle\phi_k,P\phi_k\rangle\ge1-\theta$. Since $P$ is an orthogonal
projection, $\langle\phi_k,P\phi_k\rangle=\|P\phi_k\|^2$, and Bessel's
inequality in the standard basis $e_j$ gives
\begin{align}
 \sum_{k=1}^K\|P\phi_k\|^2
 =\sum_j\sum_{k=1}^K|\langle\phi_k,Pe_j\rangle|^2
 \le\sum_j\|Pe_j\|^2=\operatorname{tr}P=n.
 \label{eq:hastings05_final_state_count_bessel}
\end{align}
Hence $(1-\theta)K\le n$. For $\theta=1/2$ this is the asserted $K\le2n$.

\section{Scope}
The argument needs orthonormal states and their energy under $H$. In the source
the $\alpha^h$ are eigenvectors of the last compressed pair operator, so they
are orthonormal. Their energy under $H$ is what the final energy estimate,
eq.~(hchy) at lines 709--720, provides after the comparison of
$\sum_iM_i$ with $H$; the count is therefore stated with the energy under $H$
as a hypothesis.
\bibliographystyle{unsrt}
\bibliography{references}
\end{document}
